\chapter{基于迁移学习的小样本电机故障诊断方法}

\section{问题分析}

在实际工业场景中，电机故障样本往往非常稀缺，难以满足深度学习模型的大样本训练需求。迁移学习通过利用源域数据上学习到的知识来帮助目标域上的学习任务，是解决小样本问题的有效途径。

\section{迁移学习策略设计}

\subsection{源域与目标域定义}

设源域$\mathcal{D}_S = \{(\mathbf{x}_i^S, y_i^S)\}_{i=1}^{N_S}$为实验室条件下采集的充足标注数据，目标域$\mathcal{D}_T = \{(\mathbf{x}_j^T, y_j^T)\}_{j=1}^{N_T}$为工业现场采集的少量标注数据，其中$N_T \ll N_S$。

\subsection{预训练-微调策略}

本文采用预训练-微调的迁移学习策略：

（1）预训练阶段：在源域数据上训练CNN模型，学习通用的振动信号特征表示。使用交叉熵损失函数：

\begin{equation}
	\mathcal{L}_{pre} = -\frac{1}{N_S}\sum_{i=1}^{N_S}\sum_{c=1}^{C} y_{i,c}^S \log \hat{y}_{i,c}^S
	\label{eq:pretrain_loss}
\end{equation}

（2）微调阶段：冻结前两层卷积块，仅微调后两层和全连接层，使用目标域数据训练：

\begin{equation}
	\mathcal{L}_{ft} = -\frac{1}{N_T}\sum_{j=1}^{N_T}\sum_{c=1}^{C} y_{j,c}^T \log \hat{y}_{j,c}^T + \lambda \|\theta\|_2^2
	\label{eq:finetune_loss}
\end{equation}

其中$\lambda$为正则化系数，$\theta$为可训练参数。

\section{实验结果与分析}

\subsection{小样本实验设置}

从CWRU数据集中随机选取不同比例的样本作为训练集，测试不同方法在小样本条件下的诊断性能。

\subsection{实验结果}

\begin{table}[htbp]
	\centering
	\caption{不同方法在小样本条件下的诊断准确率(\%)}
	\label{tab:transfer_results}
	\begin{tabular}{ccccc}
		\toprule
		训练样本数/类 & CNN & DNN & SVM & 本文方法 \\
		\midrule
		10 & 72.5 & 68.3 & 75.1 & 85.6 \\
		20 & 81.2 & 76.8 & 80.3 & 90.2 \\
		50 & 89.7 & 85.4 & 86.5 & 94.1 \\
		100 & 94.3 & 91.6 & 89.8 & 96.8 \\
		\bottomrule
	\end{tabular}
\end{table}

实验结果表明，在每类仅10个训练样本的极端小样本条件下，本文方法仍能达到85.6\%的诊断准确率，显著优于直接训练的CNN方法（72.5\%），验证了迁移学习策略在小样本故障诊断中的有效性。

\section{本章小结}

本章提出了基于迁移学习的小样本电机故障诊断方法，采用预训练-微调策略解决工业场景中小样本条件下模型训练困难的问题。实验结果表明，该方法在小样本条件下显著优于直接训练的方法。